% ------------------------------------------------------------
% Page 83
% ------------------------------------------------------------

Un petit nombre travaillait sans autorisation du STO en attendant une identité d’emprunt.
(Ceux pour lesquels on n’avait pas pu faire autrement).

\vspace{0.45cm}

D’autres s’étaient fait inscrire sous une identité d’emprunt. Il y en eut deux, dont un de
Vébron qui ne travaillèrent jamais mais dont on avait demandé l’autorisation d’embauche.
Tout cela posait un très sérieux problème administratif ; j’avais proposé, pour le résoudre, de
faire un livre de paye de circonstance pour le STO où j’inscrivais les salaires fictifs ou pas de
tous les inscrits.

\vspace{0.45cm}

Il fallu deux fois aller à Mende pour faire viser le livre que j’avais paraphé moi-même avec le
sceau de la Justice de Paix de Meyrueis que je détenais comme Greffier. C’est notre ami
Sylvain Rey, hôtelier à Meyrueis qui s’était fait embaucher comme chef de vente qui allait à
Mende faire viser le livre de paye. Il le faisait avec beaucoup d’assurance. Quelle
gymnastique de chiffres pour arriver à cadrer avec la Caisse.

\vspace{0.45cm}

Un matin de cet hiver 1944 alors que je quittais notre maison de Pradines située à 1,5 km du
village je vis deux sentinelles allemandes sur le pont de Roquedols. Ils me laissèrent passer. A
l’entrée de Meyrueis même chose. En passant devant l’Hôtel de France je vis un camion Les
allemands embarquaient les polonais. Je continuais mon chemin vers l’usine. Arrivé à la
hauteur du rocher d’Alaouze qui surplombe la route je vis tomber un petit caillou devant mon
vélo. Levant la tête j’aperçus le groupe des irréguliers dont « Séguret ». « Sont-ils partis ? »
Je fis signe que non.

\vspace{0.45cm}

Une heure après je vis arriver notre Directeur qui, me prenant à part me raconta qu’il venait
d’avoir des sueurs froides. Il avait un vieux revolver d’ordonnance dans sa chambre. Voyant
les allemands par la fenêtre il pensa que si sa propriétaire était perquisitionnée, il y aurait des
problèmes pour elle et pour lui si on trouvait cette arme. Il ficela le revolver autour de sa
jambe et sortit pour prendre la voiture de liaison de l’usine. A cet instant une voix gutturale
cria « Papiers ! ». Il sortit sa carte d’identité tandis que le revolver commençait à glisser le
long de sa jambe. Avisant l’urinoir qui était à quelques mètres il s’y engouffra et put ainsi
camoufler l’arme dans son veston.

\vspace{0.45cm}

Les polonais furent embarqués sans que nous puissions régler leur compte. Nous
n’entendîmes plus jamais parler d’eux…

\vspace{0.55cm}

\textbf{Rencontre avec le maquis}

\vspace{0.35cm}

Au début du printemps 1944, nous revenions de Nîmes Rolland, Georges, Joseph Mondino au
volant et moi-même pour les besoins de la comptabilité. Arrivés après Anduze au croisement
de la route de Lasalle, sur le pont, à quelques km de Saint Jean du Gard nous fûmes arrêtés
par le maquis. Ils étaient trois avec une moto (vraisemblablement le maquis de Lasalle).
Mains en l’air et mitraillette sur le ventre, nous fûmes fouillés : « Je te dis que ce sont des
miliciens » disait l’un. Nous avons eu beau protester, ils ont fini par nous croire et nous
relâcher. Notre chargement de charbon de bois avait été livrée comme la plupart du temps à
Nîmes à l’Agence Ford rue Hôtel Dieu. 
